\section{Introduction}\label{sec:introduction}

Real-time score following is the process of aligning a live performance with the corresponding music score as the music proceeds in real time. It is fundamental for downstream applications, including automatic accompaniment, automatic page-turning, synchronized score display, and many other interactive music systems \cite{orio2003score,dannenberg2006music}. Classical approaches use online Dynamic Time Warping (DTW) or probabilistic state models to achieve real-time alignment \cite{dixon2005line,cont2010coupled,nakamura2016real}. However, these methods require a symbolic score such as MIDI or MusicXML, which is not always available.

Image-based score following has emerged over the past decade by working directly on sheet images. Early work classified positions into coarse staff-level buckets on small local image crops \cite{dorfer2016towards}, limiting both spatial resolution and the field of view. Reinforcement-learning methods widened the view to score strips but still operated on narrow windows, making re-engagement difficult when the tracker drifted off track \cite{dorfer2018rl,henkel2019score}. Full-page segmentation addressed the field-of-view limitation by predicting a pixel-level heatmap over the entire page \cite{henkel2020learning}, but only localized at the note level without explicit system or bar tracking. The most recent of these image-based methods proposed multi-resolution detection with You Only Look Once (YOLO)-based object detection frameworks \cite{henkel2021real}. This work predicts the positions of systems, bars, and notes simultaneously using three independent output heads. However, there are no constraints or cascaded structures among these predictions. Nothing enforces that the predicted note lies within the predicted bar, or that the predicted bar lies within the predicted system. The model implicitly learns structural consistency, and disagreements among the three prediction levels are often observed.

This independence causes another problem. Because each head independently predicts over the full page, the bar and note heads must search all candidates on the page, increasing the likelihood of errors and instabilities. A cascaded structure would narrow the search at each stage: once the system is selected, only bars within that system are candidates; once a bar is selected, the candidate note positions are confined to that bar. More fundamentally, this method frames score tracking as an object detection problem. It detects from scratch at every audio frame, predicting bounding boxes as if the system and bar locations were unknown. However, in music performances, the score is static, and all candidate positions are known in advance. Detection is therefore unnecessary. The real task is to select which of the known candidates is currently being played. This formulation is feasible whenever score layout information is available. 

Additionally, existing image-based methods lack mechanisms to recover from score discontinuities such as repeats, da capo (D.C.), or coda jumps, even though such events are common in practice and have long been addressed in symbolic score following \cite{nakamura2016real,morsi2022bottlenecks,park2025matchmaker}. In our analysis of the MSMD test set, 66 out of 94 pieces contain written repeat structures, producing 131 jumps in standard performance order. Despite this prevalence, no existing real-time image-based score following method includes an explicit mechanism for handling such events during online tracking.

This paper proposes CODA, a cascaded online alignment framework that addresses these gaps. The key observation is that on a fixed score page, all system and bar locations are already known. The tracking problem is therefore reformulated as a selection task among known candidates. CODA follows the cascaded structure of music scores directly: it first selects the active system among all systems, then the active bar within that selected system, and finally the note position inside the selected bar. This formulation naturally enforces geometric consistency and narrows the search space at each stage. To handle score discontinuities, CODA models all jumps as arbitrary position changes preceded by silence, enabling a single recovery mechanism that generalizes across repeats, D.C., coda jumps, and performer errors without requiring knowledge of the score's repeat structure. The major contributions are:
\begin{itemize}
\item A cascaded selection-and-regression formulation that enforces geometric consistency across system, bar, and note predictions while narrowing the search space at each stage.
\item A causal streaming architecture combining Mamba audio encoding, FiLM conditioning, cross-attention, and beam search with learned temporal priors.
\item A silence-driven jump recovery mechanism for arbitrary score discontinuities, including repeats, D.C., and coda jumps, without requiring prior knowledge of the score's repeat structure.
\item A repeat-aware jump test benchmark with both annotated and random music discontinuities for all 94 pieces in the MSMD test set, providing the first standardized evaluation protocol for discontinuity handling in image-based score following.
\end{itemize}

